\chapter{Otitis Externa}

\section*{Definition}
Otitis externa (swimmer's ear) is an inflammation or infection of the external auditory canal, extending from the pinna to the tympanic membrane, with a duration less than 6 weeks in the acute form.

\section*{What Happens in Your Body}
Disruption of the protective ceruminous barrier and acidic pH of the ear canal allows bacterial or fungal invasion of the squamous epithelium. The resulting edema and exudate narrow the canal lumen, causing pain, itching, and conductive hearing loss.

\section*{Causes}
\begin{itemize}
  \item \textbf{Bacterial (90\%):} Pseudomonas aeruginosa and Staphylococcus aureus (most common)
  \item \textbf{Fungal (10\%):} Aspergillus niger (black spores) and Candida albicans, often after prolonged antibiotic drops
  \item \textbf{Predisposing factors:} Water exposure (swimming), high humidity, trauma (cotton swabs), hearing aids, eczema, psoriasis
  \item \textbf{Necrotizing (malignant) otitis externa:} Pseudomonas osteomyelitis of the temporal bone in diabetics and immunocompromised people
\end{itemize}

\section*{When to See a Doctor}
See your doctor if:
\begin{itemize}
  \item Severe ear pain worsened by manipulation of the pinna or tragus
  \item Purulent otorrhea with canal edema and debris
  \item Hearing loss or aural fullness
  \item Diabetes or immunocompromise with ear pain (suspect necrotizing OE)
  \item Granulation tissue in the canal floor (suspect necrotizing OE)
\end{itemize}

\section*{Related Symptoms}
\begin{itemize}
  \item Otalgia (ear pain, often severe)
  \item Pruritus (early symptom, often intense)
  \item Otorrhea (purulent or watery discharge)
  \item Conductive hearing loss
  \item Trismus (jaw pain on chewing)
\end{itemize}
\textbf{Important:} The pathognomonic sign of otitis externa is severe pain with traction on the pinna or pressure on the tragus, which is absent in uncomplicated acute otitis media.

\section*{Self-Care / Home Management}
\begin{itemize}
  \item Keep ear dry during treatment (shower cap, avoid swimming)
  \item Over-the-counter analgesic (acetaminophen or ibuprofen) for pain
  \item Avoid inserting anything into the ear canal (cotton swabs, fingers)
  \item Acetic acid 2\% drops (one dropperful tid) for mild cases without canal edema
  \item Warm compress applied to the affected ear for comfort
\end{itemize}

\section*{How Your Doctor Will Evaluate You}
\begin{itemize}
  \item Otoscopy: Erythematous, edematous canal with debris; TM appears normal (if visible)
  \item Gentle suction or wick-aided debris removal to visualize the TM
  \item Pain provocation test (tragal tenderness and pinna traction)
  \item Culture of discharge if refractory to initial empiric therapy
  \item CT temporal bone if necrotizing OE suspected (bone erosion, soft tissue extension)
\end{itemize}

\section*{Treatment Options}
\begin{itemize}
  \item Topical antibiotic drops: Aminoglycoside (gentamicin, tobramycin) or fluoroquinolone (ciprofloxacin, ofloxacin) tid for 7--10 days
  \item Combination product (neomycin/polymyxin B/hydrocortisone) if TM intact
  \item Topical fluoroquinolone monotherapy preferred if TM perforation possible (safe for middle ear)
  \item Wick placement (Pope wick or Merocel) if severe canal edema prevents drop penetration
  \item Oral antibiotics (ciprofloxacin) if infection extends beyond the canal or in immunocompromise
  \item Clotrimazole 1\% solution or thiabendazole suspension for fungal OE
  \item Necrotizing OE: IV anti-pseudomonal antibiotics (ceftazidime, ciprofloxacin) for 6--8 weeks with serial imaging
\end{itemize}

\clearpage
